% Contextual sources now cited in the supporting body argument.
% Bibliographic records have not all been independently full-text audited.
\paragraph{Interaction patterns, reasoning histories, and evaluation consequences.}
Multi-agent deliberation studies consider debate, cooperation, negotiation and exchange of intermediate messages~\cite{ext069,ext070,ext074,ext082,ext083}.
Adaptive team structures and communication-topology studies further vary who participates, which messages are shared, and how the solution process unfolds~\cite{ext084,ext085,ext086,ext087}.
These designs generate heterogeneous communication histories; a later diagnosis can require message-level provenance that aggregate task success does not preserve~\cite{ext095,ext105,ext108,ext115}.

Trajectory-level approaches cover search-structured reasoning, counterfactual rollouts and intermediate supervision~\cite{ext055,ext056,ext061}.
Agent memory surveys, self-collaboration and reflection mechanisms likewise differentiate remembered experiences from the active prompt and model state~\cite{ext092,ext093,ext094}.
We treat these as candidate sources of replay requirements and intermediate evidence, not as direct experimental results on diagnostic-retention policies~\cite{ext102,ext111,ext062}.

Long-running social simulations and goal-directed collaboration create interactions whose externally visible state may persist beyond one model request~\cite{ext088,ext089}.
Task hierarchies and institutional multi-agent representations add organizational constraints and dependencies to these histories~\cite{ext090,ext091}.
These application settings motivate evaluation beyond isolated model completions while remaining outside the core 60-item evidence census~\cite{ext109,ext143}.

